\documentclass[12pt, a4paper]{article}
\usepackage[utf8]{inputenc}
\usepackage[T1]{fontenc}
\usepackage[main=brazilian, english]{babel}
\usepackage{indentfirst}
\usepackage{geometry}



\title{Aula Metodologia da Resenha}
\author{Erickson G. Müller}
\date{\today}

\begin{document}
	\maketitle
	\section{Passo 1: }
	Fazer o cabeçalho, o nome dos autores, notinha de rodapé, referência bibliográfica NBR 6023 ABNT

	Entrar na NBR 6023, ao resenhar o artigo, dar um ctrl+F e olhar como é feita a referência bibliográfica.

	\section{Introdução}
		Objetivo: cativar o leitor. Mostrar o objetivo/problema/questão do autor. Explicar a metodologia de pesquisa.

		"Este livro, escrito por Fulano, Professor doutor da universidade Tal, trata da liderança indígena taltaltal."
	\section{Desenvolvimento}
		Resumir a estrutura do livro. "No capítulo 1 a autora fala disso aquilo". O resumo de cada capítulo é proporcional ao tamanho do capítulo. Fazer citações com frequência. As citações tem suas regras de acordo com o tamanho.
	\section{Conclusão}
		Analisar o objeto resenhado. Qualidade da abordagem, atendimento aos objetivos propostos, sustentação das teses...
\end{document}
